\documentclass[smallextended,natbib,nospthms]{svjour3}
%% Journal of Combinatorial Optimization: Springer svjour3, option smallextended.
%% Compile with pdfLaTeX. This folder must contain svjour3.cls and svglov3.clo.

\usepackage[T1]{fontenc}
\usepackage{iftex}
\ifPDFTeX
  \usepackage[utf8]{inputenc}
\fi
\usepackage{upquote}
\usepackage{microtype}
\usepackage{amsmath,amssymb,mathtools}
\usepackage{amsthm}
\usepackage{booktabs}
\usepackage{graphicx}
\usepackage{array}
\usepackage{tabularx}
\usepackage{multirow}
\usepackage{enumitem}
\usepackage{xcolor}
\usepackage{tikz}
\usetikzlibrary{arrows.meta,calc,positioning,shapes.geometric}
\usepackage{algorithm}
\usepackage[noend]{algpseudocode}
\usepackage{caption}
\usepackage{aliascnt}

% cleveref needs a dummy chapter list under svjour3
\makeatletter
\def\cl@chapter{}
\makeatother
\usepackage{hyperref}
\usepackage[nameinlink,noabbrev]{cleveref}
\emergencystretch=1.5em
\smartqed

% Make algorithmic line anchors unique across the three algorithms.
\makeatletter
\providecommand{\theHALG@line}{}
\renewcommand{\theHALG@line}{\thealgorithm.\arabic{ALG@line}}
\makeatother

\hypersetup{
  colorlinks=true,
  allcolors=black,
  pdftitle={Minimum Maximal Acyclic Matchings in Connected Cubic Graphs},
  pdfauthor={Hao Tang}
}

% Environments share the theorem counter through aliascnt so that cleveref
% reports "Lemma 4.2" instead of "Theorem 4.2".
\newtheorem{theorem}{Theorem}[section]

\newaliascnt{lemma}{theorem}
\newtheorem{lemma}[lemma]{Lemma}
\aliascntresetthe{lemma}

\newaliascnt{proposition}{theorem}
\newtheorem{proposition}[proposition]{Proposition}
\aliascntresetthe{proposition}

\newaliascnt{corollary}{theorem}
\newtheorem{corollary}[corollary]{Corollary}
\aliascntresetthe{corollary}

\theoremstyle{definition}
\newaliascnt{definition}{theorem}
\newtheorem{definition}[definition]{Definition}
\aliascntresetthe{definition}

\theoremstyle{remark}
\newaliascnt{remark}{theorem}
\newtheorem{remark}[remark]{Remark}
\aliascntresetthe{remark}

\crefname{theorem}{Theorem}{Theorems}
\Crefname{theorem}{Theorem}{Theorems}
\crefname{lemma}{Lemma}{Lemmas}
\Crefname{lemma}{Lemma}{Lemmas}
\crefname{proposition}{Proposition}{Propositions}
\Crefname{proposition}{Proposition}{Propositions}
\crefname{corollary}{Corollary}{Corollaries}
\Crefname{corollary}{Corollary}{Corollaries}
\crefname{definition}{Definition}{Definitions}
\Crefname{definition}{Definition}{Definitions}
\crefname{remark}{Remark}{Remarks}
\Crefname{remark}{Remark}{Remarks}

\newcommand{\acnum}{\mu'_{\mathrm{ac}}}
\newcommand{\maxacnum}{\nu_{\mathrm{ac}}}
\newcommand{\LowAcy}{\textnormal{\textsc{Low-Acy-Matching}}}
\newcommand{\ceil}[1]{\left\lceil #1\right\rceil}
\newcommand{\floor}[1]{\left\lfloor #1\right\rfloor}
\newcommand{\bigO}{\mathcal{O}}
\DeclareMathOperator{\Forest}{Forest}
\DeclareMathOperator{\popcount}{popcount}

\algrenewcommand\algorithmicrequire{\textbf{Input:}}
\algrenewcommand\algorithmicensure{\textbf{Output:}}

\tikzset{
  vertex/.style={circle,draw,fill=white,inner sep=0pt,minimum size=5.2mm,font=\scriptsize},
  fcomp/.style={rectangle,draw,align=center,minimum width=14mm,minimum height=8mm,font=\scriptsize},
  htree/.style={circle,draw,align=center,minimum size=11mm,font=\scriptsize},
  hcycle/.style={regular polygon,regular polygon sides=3,draw,align=center,
                 minimum size=12mm,font=\scriptsize}
}

\journalname{Journal of Combinatorial Optimization}

\begin{document}

\title{Minimum Maximal Acyclic Matchings in Connected Cubic Graphs:
Exhaustive Enumeration and a Sharp Bound}
\titlerunning{Minimum maximal acyclic matchings in cubic graphs}
\author{Hao Tang}
\authorrunning{H. Tang}
\institute{H.~Tang \at
College of Software, Jilin University, Changchun 130012, China \\
\email{tanghao5524@mails.jlu.edu.cn}}
\date{Received: date / Accepted: date}

\maketitle

\begin{abstract}
An acyclic matching is a matching whose endpoints induce a forest.
Given a graph $G$, the problem \LowAcy{} asks for an inclusion-wise maximal
acyclic matching of minimum cardinality; we write $\acnum(G)$ for this
optimum. The decision version is NP-complete on general graphs. We determine
the exact extremal value of $\acnum$ over connected simple cubic graphs.

F\"urst and Rautenbach proved that the maximum acyclic matching number
satisfies $\maxacnum(G)\geq n/4$ on connected cubic graphs other than
$K_{3,3}$. The corresponding $\floor{n/4}$ benchmark holds for $\acnum$
through order 14, but fails at order 16: four of $4060$ non-isomorphic graphs
have $\acnum(G)=3$. For a maximal acyclic matching $M$, analysis of the
unmatched induced subgraph and its incidence with the forest induced by
$V(M)$ yields $n\leq 5|M|$ when $|M|$ is even and $n\leq 5|M|+1$ when $|M|$
is odd. Explicit connected cubic constructions attain the bound at every even
order. Thus, if $f_3(n)$ denotes the minimum of $\acnum(G)$ over connected
simple cubic graphs of order $n$, then
$f_3(n)=\ceil{(n-1)/5}$ for every even $n\geq 4$.

Two independent exact solvers, dual regular-graph generators, and an
endpoint-set audit certify the census through order 16. As a direct counting
consequence, deciding $\acnum(G)\leq K$ on this class reduces immediately to
instances with at most $5K+1$ vertices, and is therefore fixed-parameter
tractable in the solution-size parameter $K$.
\keywords{Acyclic matching \and Cubic graph \and Extremal graph theory
\and Exhaustive enumeration \and Parameterized algorithm}
\end{abstract}

\section{Introduction}

Matching is a fundamental problem in combinatorial optimization and graph
algorithms. Classical maximum matching is polynomial-time solvable
\cite{Edmonds1965,LovaszPlummer2009}, whereas the complexity often increases
sharply when feasibility is defined through an induced-subgraph property or
when the objective is changed from a maximum feasible solution to a minimum
maximal one. Let $G$ be a finite simple graph and let $M\subseteq E(G)$ be a
matching. The matching $M$ is \emph{acyclic} if its endpoints induce a forest
in $G$. Acyclic matchings were developed systematically within the framework
of generalized subgraph-restricted matchings \cite{Goddard2005}, and their
maximum and minimum-maximal variants have since led to distinct algorithmic
lines of research
\cite{Chaudhary2024,Chaudhary2025,Furst2019,Panda2023,Zehavi2025}.

This paper studies \LowAcy: among all inclusion-wise maximal acyclic matchings,
find one of minimum cardinality. Its optimum is denoted by $\acnum(G)$.
This problem differs from maximum acyclic matching and from ordinary minimum
maximal matching, which is closely related to edge domination
\cite{Yannakakis1980}. Induced and uniquely restricted matchings provide
further nearby reference points \cite{Cameron1989,Golumbic2001}, and the
minimum-maximal variant of uniquely restricted matching exhibits the same
min--max pattern studied here \cite{ChaudharyPanda2021}.
The decision version of \LowAcy{} is NP-complete on general graphs. For
restricted classes, it remains NP-complete on bipartite graphs of maximum
degree six and is APX-hard on 4-regular graphs, whereas proper interval graphs
admit a linear-time algorithm \cite{Chaudhary2024,Chaudhary2025}. These
results do not settle the per-instance complexity on connected cubic graphs.
The present paper addresses first a different, extremal question: the minimum
possible value of $\acnum(G)$ among connected cubic graphs of a fixed order.
\citet{Chaudhary2024} observed that even a general lower
bound was missing for fixed $r$-regular graphs with $r\geq 3$. Cubic graphs
are therefore the first unresolved regular class in which to seek a sharp
universal bound.

The corresponding question for the \emph{maximum} acyclic matching number
$\maxacnum$ is settled. \citet{FurstRautenbach2018}
proved that every subcubic graph without isolated vertices satisfies
\[
  \maxacnum(G)\geq
  \tfrac14\bigl(n-\kappa_G(K_{2,3})-\kappa_G(K_4^{+})-2\kappa_G(K_{3,3})\bigr),
\]
where $\kappa_G(H)$ counts the components of $G$ isomorphic to $H$ and
$K_4^{+}$ arises from $K_4$ by subdividing one edge once. Neither $K_{2,3}$
nor $K_4^{+}$ is cubic, so the bound specialises to $\maxacnum(G)\geq n/4$ for
every connected cubic graph other than $K_{3,3}$, and it is tight. Since
$\acnum(G)\leq\maxacnum(G)$, this inequality points in the wrong direction for
the minimum-maximal parameter and settles nothing about $\acnum$. It does,
however, make $\floor{n/4}$ the natural first guess for $f_3(n)$, and our
census shows that the guess is wrong: the quarter threshold that governs the
maximum variant on cubic graphs does not govern the minimum-maximal one, whose
true growth rate is $n/5$.

We follow an exact-computation-to-structure route: discover counterexamples
computationally, formulate a structural conjecture, and then determine the
extremal function by an independent proof. We first implement a
definition-level solver and an optimized bitset solver, and then enumerate
small non-isomorphic connected cubic graphs. The census through order 14
happens to satisfy the benchmark $\acnum(G)\geq\floor{n/4}$ inherited from
$\maxacnum$, but the
order-16 census refutes it and identifies four smallest counterexamples.
Their structure shows that the decisive feature is not the local form of the
matching edges, but how every edge in the subgraph induced by unmatched
vertices is blocked by connected components of the forest induced by the
matching endpoints. This observation yields a local characterization, a
component classification, and an exact cut-edge count.

\subsection{Main results and computational validation}

The contributions are as follows.
\begin{itemize}[leftmargin=1.5em]
  \item For every maximal acyclic matching $M$ with $k=|M|$, we prove
  $n\leq 5k$ when $k$ is even and $n\leq 5k+1$ when $k$ is odd. The split
  identifies precisely when the one-vertex slack can occur.

  \item Equality constructions are given for every even $n\geq 4$, determining
  the exact extremal function
  \[
    \begin{aligned}
      f_3(n)
      &=\min\bigl\{\acnum(G):
        G\text{ connected simple cubic of order }n\bigr\}\\
      &=\ceil{\frac{n-1}{5}}.
    \end{aligned}
  \]

  \item All non-isomorphic connected simple cubic graphs of orders
  $n=4,6,\ldots,16$ are evaluated exactly. The complete histogram exposes
  four of the 4060 order-16 graphs with $\acnum(G)=3<4$ and leads to the sharp
  theorem. Exactly three of these four are realised by the equality
  construction of \Cref{sec:constructions}; the fourth is not, so extremal
  graphs are not confined to a single structural family.

  \item For reproducibility, a definition-level edge-subset solver and a bitset
  backtracking solver are specified and proved correct. Their results are
  cross-checked through order 12, while a structurally independent
  endpoint-set solver certifies every optimum at orders 14 and 16. GENREG and
  nauty/\texttt{geng} produce identical canonical input censuses.

  \item As a post-theorem consequence, we derive an exact solver that contracts
  selected matching edges, maintains an edge-interaction forest with rollback
  union--find, and uses a theorem-derived defect budget. The same counting
  bound implies that deciding $\acnum(G)\leq K$ is fixed-parameter tractable
  in $K$ on connected cubic graphs: every yes-instance has at most $5K+1$
  vertices, so larger instances are immediate no-instances.
\end{itemize}

The last item is deliberately downstream of the theorem. The new solver uses
the proved lower bound and therefore is not counted as independent evidence
for that bound.

\subsection{Organization}

\Cref{sec:definitions} gives the definitions, two theorem-independent exact
algorithms, and their correctness analyses. \Cref{sec:census} reports the
non-isomorphic census and counterexample certificates.
\Cref{sec:blocking} establishes the local blocking structure, and
\Cref{sec:bound} proves the sharp lower bound. \Cref{sec:constructions} gives
equality constructions for all even orders. \Cref{sec:solverc} develops the
post-theorem interaction-forest solver and records the FPT consequence of the
bound. \Cref{sec:verification} describes independent verification,
reproducibility, and threats to validity.
\Cref{sec:discussion,sec:conclusion} discuss open questions and conclude the
paper. Detailed Solver~C timings, equality-family scaling, and ablations are
collected in \Cref{app:benchmark}.

\section{Problem definition and exact algorithms}
\label{sec:definitions}

\subsection{Definitions and decision semantics}

All graphs considered in this paper are finite, undirected, and simple. For
$X\subseteq V(G)$, let $G[X]$ denote the vertex-induced subgraph of $G$ on
$X$. For a matching $M$, let $V(M)$ be the set of all endpoints of edges in
$M$.

\begin{definition}
A matching $M$ is \emph{acyclic} if $G[V(M)]$ is a forest. It is
\emph{maximal acyclic} if it is not strictly contained in any larger acyclic
matching. Define
\[
  \acnum(G)=\min\{|M|:M\text{ is a maximal acyclic matching of }G\}.
\]
For every even integer $n\geq 4$, define
\[
  f_3(n)=
  \min\{\acnum(G):G\text{ is a connected simple cubic graph of order }n\}.
\]
\end{definition}

A common implementation error is to interpret an extension as merely adding a
new edge $uv$ to $G[V(M)]$. After $uv$ is added to the matching, however, one
must test the full vertex-induced subgraph
$G[V(M)\cup\{u,v\}]$, including every edge from $u$ and $v$ to the previous
matching endpoints. All exact solvers in this work implement this semantics.

\subsection{Solver A: definition-level enumeration}

\begin{algorithm}[t]
\caption{Solver A: definition-level exact enumeration}
\label{alg:solver-a}
\begin{algorithmic}[1]
\Require A finite simple graph $G=(V,E)$
\For{$k=0,1,\ldots,|E|$}
  \ForAll{$k$-edge subsets $M\subseteq E$}
    \If{$M$ is not a matching}
      \State \textbf{continue}
    \EndIf
    \If{$G[V(M)]$ is not a forest}
      \State \textbf{continue}
    \EndIf
    \State $U\gets V\setminus V(M)$; \quad $\mathit{maximal}\gets\textsc{True}$
    \ForAll{$uv\in E(G[U])$}
      \If{$G[V(M)\cup\{u,v\}]$ is a forest}
        \State $\mathit{maximal}\gets\textsc{False}$; \quad \textbf{break}
      \EndIf
    \EndFor
    \If{$\mathit{maximal}$}
      \State \Return $(k,M)$
    \EndIf
  \EndFor
\EndFor
\end{algorithmic}
\end{algorithm}

Solver A exhausts every candidate of smaller cardinality before advancing to
the next cardinality, so the first returned solution is optimal. Deliberately
avoiding sophisticated pruning, it serves as a definition-level baseline and
as an independent cross-check. If $m=|E(G)|$, its worst-case running time is
$\bigO(2^m\operatorname{poly}(n+m))$, so it is suitable only for small graphs.

\begin{proposition}
\label{prop:solver-a}
Solver A returns a minimum-cardinality maximal acyclic matching.
\end{proposition}

\begin{proof}
The algorithm returns only if $M$ is a matching, $G[V(M)]$ is a forest, and no
edge $uv\in E(G[V\setminus V(M)])$ extends it while preserving the induced
forest condition. Thus the returned solution is feasible and maximal. The
outer loop increases $k$ only after all smaller edge subsets have been checked,
so no smaller solution exists.
\end{proof}

\subsection{Solver B: bitset increasing-cardinality search}

Index the edges as $E=(e_0,\ldots,e_{m-1})$ and let $p[i]$ be the bit mask of
the two endpoints of $e_i$. Solver B searches in increasing target size and
caches the exact predicate $\Forest(X)$, which tests whether $G[X]$ is a
forest. For a forest $G[X]$, the cache stores both acyclicity and its component
count. Cyclicity is monotone under adding vertices to an induced subgraph, so
a failed forest test safely prunes the entire branch. Two capacity tests are
also safe: a completion needs at least \textit{need} eligible suffix edges and
$2\,\textit{need}$ distinct endpoints.

\begin{algorithm}[t]
\caption{Solver B: bitset- and cache-driven exact search}
\label{alg:solver-b}
\begin{algorithmic}[1]
\Require $G=(V,E)$ with endpoint masks $p[0],\ldots,p[m-1]$
\Function{Maximal}{$X$}
  \For{$i=0,\ldots,m-1$}
    \If{$(p[i]\mathbin{\&}X)=0$ and $\Forest(X\mathbin{|}p[i])$}
      \State \Return \textsc{False}
    \EndIf
  \EndFor
  \State \Return \textsc{True}
\EndFunction
\Function{Search}{$\mathit{start},q,X,L,k$}
  \If{$q=k$}
    \State \Return $L$ if \Call{Maximal}{$X$}; otherwise failure
  \EndIf
  \State $\mathit{need}\gets k-q$
  \State $R\gets\{i\geq\mathit{start}:(p[i]\mathbin{\&}X)=0\}$
  \If{$|R|<\mathit{need}$}
    \State \Return failure
  \EndIf
  \State $Y\gets\bigvee_{i\in R}p[i]$
  \If{$\popcount(Y)<2\,\mathit{need}$}
    \State \Return failure
  \EndIf
  \ForAll{$i\in R$ in increasing order}
    \State $X'\gets X\mathbin{|}p[i]$
    \If{$\neg\Forest(X')$}
      \State \textbf{continue}
    \EndIf
    \State $A\gets\Call{Search}{i+1,q+1,X',L\cup\{e_i\},k}$
    \If{$A\neq\text{failure}$}
      \State \Return $A$
    \EndIf
  \EndFor
  \State \Return failure
\EndFunction
\For{$k=0,1,\ldots,\floor{n/2}$}
  \State $A\gets\Call{Search}{0,0,0,\varnothing,k}$
  \If{$A\neq\text{failure}$}
    \State \Return $(k,A)$
  \EndIf
\EndFor
\end{algorithmic}
\end{algorithm}

\begin{proposition}
\label{prop:solver-b}
Solver B returns a minimum-cardinality maximal acyclic matching.
\end{proposition}

\begin{proof}
At every recursive call, $L$ is an acyclic matching, $X=V(L)$, and the edge
indices in $L$ are strictly smaller than \textit{start}. An endpoint conflict
cannot occur in a matching. If too few eligible suffix edges or too few
distinct suffix endpoints remain, no completion exists. If $G[X']$ contains
a cycle, every induced supergraph on a superset of $X'$ contains that cycle.
Thus all pruning rules are safe.

Conversely, list the edge indices of any acyclic $k$-edge matching in increasing
order. The recursion can follow these indices. None conflicts with earlier
endpoints, none fails the exact forest test, and the unused edges of the
matching witness both capacity requirements. Hence every feasible $k$-edge
matching is reached exactly once.

At depth $k$, \Call{Maximal}{$X$} tests every edge with both endpoints outside
$X$. It returns true exactly when every one-edge induced extension contains a
cycle. This is equivalent to inclusion-wise maximality because acyclicity is
hereditary. The outer loop exhausts all smaller cardinalities first.
\end{proof}

The worst-case search remains exponential. Solver B contributes an exact,
auditable implementation with strong practical pruning, not a new worst-case
exponential base.

\section{Exhaustive enumeration of non-isomorphic cubic graphs}
\label{sec:census}

\subsection{Graph generation and experimental protocol}

A cubic graph has even order. The experiments process
$n=4,6,8,10,12,14,16$. To remove dependence on one generation algorithm,
every order was generated independently by GENREG, which implements
Meringer's orderly regular-graph generation method \cite{Meringer1999}, and by
nauty 2.9.3 \texttt{geng} with the restrictions
\texttt{-c -d3 -D3} and exactly $3n/2$ edges \cite{McKayPiperno2014}.
GENREG short-code records were decoded by a separate parser. Both collections
were converted to \texttt{graph6} and rechecked for simplicity, connectivity,
and degree three.

Each collection was canonically labelled with nauty \texttt{labelg}, sorted,
and compared as a multiset. At orders
$4,6,8,10,12,14,16$, the two generators produced
$1,2,5,19,85,509,4060$ distinct graphs, respectively; the canonical multisets
were identical at every order, with no missing, additional, or duplicated
class. These well-known census counts are completeness checks rather than new
enumeration results. Canonicalization is used only for comparison and is not
part of either generation procedure.

The solver protocol has three levels. First, definition-level unit tests cover
$K_4$, $K_{3,3}$, paths, cycles, and small graph-atlas instances. Second,
Solvers A and B are compared graph by graph on every cubic graph of order at
most 12. Third, an independently written endpoint-set solver computes the
optimum of every one of the 509 order-14 and 4060 order-16 graphs. It
enumerates vertex sets $S$ in increasing even cardinality, uses a disjoint-set
union test for $G[S]$, independently searches for a perfect matching of
$G[S]$, and directly tests every two-vertex induced extension. The experiment
terminates immediately on any discrepancy.

\subsection{Complete census}

\begin{table}[t]
\centering
\caption{Exact census of connected simple cubic graphs through order 16. The
benchmark column is $\floor{n/4}$.}
\label{tab:census}
\begin{tabular}{rrrrrr}
\toprule
$n$ & Classes & Min. & Max. & Benchmark & At min. / violations\\
\midrule
4  & 1    & 1 & 1 & 1 & 1 / 0\\
6  & 2    & 1 & 1 & 1 & 2 / 0\\
8  & 5    & 2 & 2 & 2 & 5 / 0\\
10 & 19   & 2 & 3 & 2 & 14 / 0\\
12 & 85   & 3 & 3 & 3 & 85 / 0\\
14 & 509  & 3 & 4 & 3 & 418 / 0\\
16 & 4060 & 3 & 4 & 4 & 4 / 4\\
\bottomrule
\end{tabular}
\end{table}

\begin{table}[t]
\centering
\caption{Complete histogram of $\acnum$ over the cubic census.}
\label{tab:histogram}
\begin{tabular}{rrr@{\qquad}rrr}
\toprule
$n$ & $\acnum$ & Classes & $n$ & $\acnum$ & Classes\\
\midrule
4  & 1 & 1  & 12 & 3 & 85\\
6  & 1 & 2  & 14 & 3 & 418\\
8  & 2 & 5  & 14 & 4 & 91\\
10 & 2 & 14 & 16 & 3 & 4\\
10 & 3 & 5  & 16 & 4 & 4056\\
\bottomrule
\end{tabular}
\end{table}

\Cref{tab:census} summarises the outcome and \Cref{tab:histogram} gives the
full histogram. For $n\leq 14$, every graph satisfies the
$\floor{n/4}$ benchmark. The first violations occur
at order 16: 4056 of the 4060 graphs have parameter value four, while four
have value three. Hence $f_3(16)=3<4=\floor{16/4}$.

\subsection{The four order-16 counterexamples}

The \texttt{graph6} encoding of the first counterexample is
\verb|O^o?GGB?w????B@@`@??X|.
Take
\[
  M=\{(0,4),(5,6),(10,12)\}.
\]
The subgraph induced by $V(M)$ is exactly $3K_2$ and is therefore a forest.
The unmatched induced subgraph contains nine edges; for every such edge
$uv$, the graph $G[V(M)\cup\{u,v\}]$ contains a triangle or a 4-cycle.
Thus $M$ cannot be extended. Solver A exhausts cardinalities zero, one, and
two, and Solver B returns the same optimum. The independent endpoint-set
audit also returns three for exactly the four canonical classes listed in
\Cref{tab:counterexamples}, whose \texttt{graph6} encodings are recorded in
\Cref{tab:graph6}. \Cref{fig:blocking} illustrates the three blocking
mechanisms established below.

\begin{table}[t]
\centering
\caption{All four counterexamples of order 16. ``Diam.'', ``Br.'',
``Tri.'', and $C_4$ denote diameter, number of bridges, number of triangles,
and number of 4-cycles, respectively.}
\label{tab:counterexamples}
\begin{tabular}{clrrrr}
\toprule
ID & Optimal matching & Diam. & Br. & Tri. & $C_4$\\
\midrule
4055 & $\{(10,12),(0,4),(5,6)\}$  & 9 & 2 & 6 & 7\\
4056 & $\{(0,4),(5,6),(10,11)\}$  & 9 & 2 & 4 & 11\\
4058 & $\{(0,4),(5,6),(11,15)\}$  & 8 & 2 & 5 & 8\\
4059 & $\{(0,4),(5,6),(11,12)\}$  & 6 & 3 & 6 & 9\\
\bottomrule
\end{tabular}
\end{table}

\begin{table}[t]
\centering
\caption{\texttt{graph6} encodings of the four counterexamples.}
\label{tab:graph6}
\begin{tabular}{rl}
\toprule
ID & \texttt{graph6} encoding\\
\midrule
4055 & \verb|O^o?GGB?w????B@@`@??X|\\
4056 & \verb|O^o?GGB?w??@?A@@`@??L|\\
4058 & \verb|O^o?GGB?w??@?B@?`@??J|\\
4059 & \verb|O^o?GGB?wW?@?@?@??W?F|\\
\bottomrule
\end{tabular}
\end{table}

\begin{figure}[t]
\centering
\begin{tikzpicture}[x=1.05cm,y=0.78cm]
  \foreach \i/\x/\y in {
    3/-3.0/0.0,1/-2.3/0.55,2/-1.4/0.0,0/-0.7/0.55,4/-0.7/1.65,
    14/0.15/2.55,15/0.25/3.45,10/1.35/3.20,11/1.15/4.25,
    12/2.25/4.10,13/2.35/5.10,5/2.55/6.05,7/2.45/7.00,
    9/3.40/7.45,6/4.10/6.95,8/4.85/7.55}
    \node[vertex] (v\i) at (\x,\y) {\i};

  \foreach \a/\b in {
    0/2,0/3,1/2,1/3,1/4,2/3,4/14,
    5/7,5/13,6/8,6/9,7/8,7/9,8/9,
    10/14,10/15,11/12,11/13,11/15,12/13,14/15}
    \draw[densely dashed,gray!75] (v\a)--(v\b);
  \draw[very thick] (v0)--(v4);
  \draw[very thick] (v5)--(v6);
  \draw[very thick] (v10)--(v12);
\end{tikzpicture}
\caption{A 16-vertex cubic graph with $\acnum(G)=3<4$. Thick edges form the
matching; dashed edges are the remaining graph edges.}
\label{fig:n16}
\end{figure}

The four counterexamples are pairwise non-isomorphic, non-bipartite, contain
triangles, and have vertex connectivity one. Their common structure motivates
the analysis of the components of the unmatched induced subgraph. In all four,
the optimal matching satisfies $c=3$ and $t=1$ in the notation of
\Cref{sec:blocking}, and the induced forest $G[V(M)]$ is $3K_2$; the
unmatched induced subgraph is $2C_3\cup P_4$ for IDs 4055, 4056, and 4058, and
$3C_3\cup K_1$ for ID 4059.

This split is exactly the boundary of the equality family. The existence
proof in \Cref{sec:constructions} uses one canonical assignment of degree
slots, recorded there. Independently enumerating all bijections of the
$t=1$ base construction yields $3456$ labelled graphs falling into precisely
three isomorphism classes, namely IDs 4055, 4056, and 4058. The graph in
\Cref{fig:n16} is one of them. ID 4059 is not obtained by any slot
assignment: it is the graph formed by three leaf pieces $K_2+C_3$ whose three
$K_2$ components each send one degree slot to a single common unmatched
vertex. Its unmatched induced subgraph therefore has an isolated vertex as
its unique tree component instead of the root $P_4$, and it has three bridges
rather than two. Consequently the construction of
\Cref{sec:constructions} is sufficient for sharpness but is not a
classification of extremal graphs; see \Cref{rem:non-uniqueness}.

\section{Local blocking structure of maximality}
\label{sec:blocking}

Fix a connected simple cubic graph $G$ and an arbitrary maximal acyclic
matching $M$. Let
\[
  k=|M|,\qquad S=V(M),\qquad U=V(G)\setminus S,\qquad
  F=G[S],\qquad H=G[U].
\]
By definition, $F$ is a forest and $|S|=2k$. Note that $k\geq 1$: the empty
matching is acyclic but never maximal, because a single edge always induces a
forest. Hence $c\geq 1$ throughout, and $U\neq\varnothing$, since otherwise
$F=G$ would be a cubic forest.

\subsection{Maximality as non-extendability by one edge}

\begin{lemma}
\label{lem:one-edge}
An acyclic matching $M$ is maximal if and only if there is no edge
$e\in E(G)\setminus M$ such that $M\cup\{e\}$ is an acyclic matching.
\end{lemma}

\begin{proof}
If an acyclic matching $M'$ strictly contains $M$, choose
$e\in M'\setminus M$. Acyclic matchings are closed under taking subsets, so
$M\cup\{e\}$ remains acyclic. Hence the absence of a one-edge extension is
equivalent to the absence of every strict extension.
\end{proof}

An edge can be added to $M$ exactly when both endpoints lie in $U$; the
candidate extension edges are precisely $E(H)$. Consequently, for every
$uv\in E(H)$, the induced subgraph $G[S\cup\{u,v\}]$ contains a cycle.

\subsection{An exact local blocking criterion}

For $x\in U$ and a component $C$ of $F$, define
\[
  d_C(x)=|N_G(x)\cap V(C)|.
\]

\begin{lemma}
\label{lem:blocking}
Let $uv\in E(H)$. Then $G[S\cup\{u,v\}]$ contains a cycle if and only if at
least one of the following holds:
\begin{enumerate}[label=(\roman*),leftmargin=2em]
  \item $d_C(u)\geq 2$ for some component $C$ of $F$;
  \item $d_C(v)\geq 2$ for some component $C$ of $F$;
  \item $u$ and $v$ are both adjacent to the same component $C$ of $F$.
\end{enumerate}
\end{lemma}

\begin{proof}
If (i) holds, two distinct neighbors of $u$ in the tree $C$ are joined by a
unique path in $C$; that path together with $u$ forms a cycle. The proof for
(ii) is symmetric. If (iii) holds, choose
$a\in N(u)\cap V(C)$ and $b\in N(v)\cap V(C)$. If $a=b$, then $uav$ is a
triangle. Otherwise the unique $a$--$b$ path in $C$, together with
$ua,uv,vb$, is a cycle.

Conversely, let $Q$ be a cycle in $G[S\cup\{u,v\}]$. Since $F$ is a forest,
$Q$ contains $u$ or $v$. If it contains exactly one, that vertex has two
neighbors on $Q$ in one component of $F$, giving (i) or (ii). If $Q$ contains
both $u$ and $v$, either deleting $uv$ from $Q$, or taking either $u$--$v$ arc
when $uv\notin E(Q)$, gives a path whose internal vertices lie in one
$F$-component. Thus (iii) holds.
\end{proof}

\begin{figure}[t]
\centering
\begin{tikzpicture}[scale=0.72, every node/.style={font=\scriptsize}]
  \begin{scope}[xshift=0cm]
    \node[vertex] (u1) at (0,2) {$u$};
    \node[vertex] (v1) at (2,2) {$v$};
    \node[vertex] (a1) at (-0.7,0) {};
    \node[vertex] (b1) at (0.2,-0.4) {};
    \node[vertex] (c1) at (1.0,0) {};
    \draw[very thick] (u1)--(v1);
    \draw[double,double distance=1.8pt,line width=0.45pt] (a1)--(b1)--(c1);
    \draw (u1)--(a1) (u1)--(c1);
    \node at (0.6,-1.0) {$u$ self-blocks};
  \end{scope}
  \begin{scope}[xshift=5.2cm]
    \node[vertex] (u2) at (0,2) {$u$};
    \node[vertex] (v2) at (2,2) {$v$};
    \node[vertex] (a2) at (1.0,0) {};
    \node[vertex] (b2) at (1.8,-0.4) {};
    \node[vertex] (c2) at (2.7,0) {};
    \draw[very thick] (u2)--(v2);
    \draw[double,double distance=1.8pt,line width=0.45pt] (a2)--(b2)--(c2);
    \draw (v2)--(a2) (v2)--(c2);
    \node at (1.8,-1.0) {$v$ self-blocks};
  \end{scope}
  \begin{scope}[xshift=10.8cm]
    \node[vertex] (u3) at (0,2) {$u$};
    \node[vertex] (v3) at (2.6,2) {$v$};
    \node[vertex] (a3) at (0.2,0) {};
    \node[vertex] (b3) at (1.2,-0.4) {};
    \node[vertex] (c3) at (2.4,0) {};
    \draw[very thick] (u3)--(v3);
    \draw[double,double distance=1.8pt,line width=0.45pt] (a3)--(b3)--(c3);
    \draw (u3)--(a3) (v3)--(c3);
    \node at (1.3,-1.0) {$u,v$ meet one $F$-component};
  \end{scope}
\end{tikzpicture}
\caption{The three local mechanisms that block an unmatched edge. Double
strokes mark paths that lie in a component of the forest $F$; these remain
legible in greyscale.}
\label{fig:blocking}
\end{figure}

\subsection{Components of the unmatched induced subgraph}

\begin{lemma}
\label{lem:h-components}
Every connected component of $H=G[U]$ is a path, a cycle, or $K_{1,3}$.
\end{lemma}

\begin{proof}
Suppose $x\in U$ has degree three in $H$. Since $G$ is cubic, $x$ has no
neighbor in $S$. For any $y\in N_H(x)$, apply \Cref{lem:blocking} to $xy$.
The vertex $x$ cannot block the edge by itself, and $x,y$ cannot be jointly
adjacent to an $F$-component. Hence $y$ must block the edge by itself. It has
exactly two neighbors in one $F$-component and, by cubicity, $x$ is its only
neighbor in $H$. Thus the component containing $x$ is $K_{1,3}$.

If an $H$-component contains no vertex of degree three, its maximum degree is
at most two. Every finite connected graph of maximum degree at most two is a
path or a cycle.
\end{proof}

\begin{lemma}
\label{lem:component-incidence}
Every tree component of $H$ is adjacent to at most three distinct components
of $F$, whereas every cycle component of $H$ is adjacent to exactly one
component of $F$.
\end{lemma}

\begin{proof}
An isolated vertex of $H$ has three neighbors in $S$ and hence meets at most
three $F$-components. If the component is $K_{1,3}$, its center has no
neighbor in $S$. Applying \Cref{lem:blocking} to each center--leaf edge shows
that the two $S$-neighbors of each leaf lie in one $F$-component; the three
leaves therefore meet at most three such components.

Consider a path $x_1\ldots x_r$. The case $r=2$ is immediate: if one endpoint
self-blocks, it contributes one $F$-component and the other at most two; if
neither self-blocks, their two sets of incident components intersect, so their
union has size at most three. For $r\geq 3$, every internal vertex has exactly
one neighbor in $S$ and cannot self-block. Applying \Cref{lem:blocking} along
the edges between consecutive internal vertices shows that all internal
vertices meet one common $F$-component $C$. Each endpoint contributes at most
one component outside $C$, so the whole path meets at most three components.
The same argument also covers $P_3$.

Finally, every vertex in a cycle component of $H$ has exactly one neighbor in
$S$ and cannot self-block. Each cycle edge must therefore be blocked through
a common $F$-component. Propagating around the cycle shows that all its
vertices meet one and the same component of $F$.
\end{proof}

\subsection{The component-incidence graph}

Let $c$ be the number of connected components of $F$, and let $t$ be the number
of tree components of $H$, including paths, isolated vertices, and
$K_{1,3}$ components. Construct a simple bipartite graph $B$ whose left side
consists of the $F$-components and whose right side consists of the
$H$-components. Two component vertices are adjacent when at least one edge
of $G$ joins the corresponding components.

\begin{lemma}
\label{lem:incidence-bound}
The inequalities $c\leq 2t+1$ and $c\leq k$ hold.
\end{lemma}

\begin{proof}
Connectivity of $G$ implies that $B$ is connected. By
\Cref{lem:component-incidence}, every cycle component of $H$ is a leaf of $B$.
Delete all such leaves to obtain $B_0$. Deleting leaves cannot disconnect the
remaining vertices. Hence $B_0$ is connected.

If $t=0$, connectedness forces $c=1$, so $c\leq 2t+1$. If $t\geq 1$, then
$B_0$ has $c+t$ vertices and at least $c+t-1$ edges. Each of its $t$
right-side vertices has degree at most three by
\Cref{lem:component-incidence}, so $|E(B_0)|\leq 3t$. Therefore
$c+t-1\leq 3t$, or $c\leq 2t+1$.

Every component of $F$ contains at least one edge of $M$, and distinct
$F$-components contain distinct matching edges. Thus $c\leq k$.
\end{proof}

\section{Sharp lower bound}
\label{sec:bound}

\subsection{Cut-edge counting}

The forest $F$ has $2k$ vertices and $c$ components, so
$|E(F)|=2k-c$. Counting the cut $E(S,U)$ from the $S$ side gives
\begin{equation}
  |E(S,U)|
  =3|S|-2|E(F)|
  =6k-2(2k-c)
  =2k+2c.
  \label{eq:cut-s}
\end{equation}
By \Cref{lem:h-components}, every component of $H$ is a tree or a cycle. A
tree component has one fewer edge than vertices, while a cycle component has
as many edges as vertices. Since $H$ has $t$ tree components,
$|E(H)|=|U|-t$. Counting the same cut from the $U$ side gives
\begin{equation}
  |E(S,U)|
  =3|U|-2|E(H)|
  =3|U|-2(|U|-t)
  =|U|+2t.
  \label{eq:cut-u}
\end{equation}
Equating \eqref{eq:cut-s} and \eqref{eq:cut-u} yields
\begin{equation}
  |U|=2k+2c-2t,
  \qquad
  n=|S|+|U|=4k+2(c-t).
  \label{eq:master-count}
\end{equation}

\subsection{Integer optimization}

Subtracting $t$ from the bounds in \Cref{lem:incidence-bound} gives
\[
  c-t\leq\min\{k-t,t+1\}.
\]
For integral $t\geq 0$, the maximum of the right-hand side is
$\ceil{k/2}$. Substituting this into \eqref{eq:master-count} gives
\[
  n\leq 4k+2\ceil{\frac{k}{2}}
  =
  \begin{cases}
    5k,   & k\text{ even},\\
    5k+1, & k\text{ odd}.
  \end{cases}
\]

\begin{theorem}
\label{thm:sharp-lower}
Let $G$ be a connected simple cubic graph of order $n$, and let $M$ be any
maximal acyclic matching of $G$ with $k=|M|$. Then $n\leq 5k$ when $k$ is
even and $n\leq 5k+1$ when $k$ is odd. In particular,
\[
  \acnum(G)\geq\ceil{\frac{n-1}{5}}.
\]
\end{theorem}

\begin{proof}
The parity-sensitive inequalities follow from \eqref{eq:master-count} and the
integer optimization above. Both imply $n\leq 5k+1$, and hence
$k\geq(n-1)/5$. Since $k$ is integral,
$k\geq\ceil{(n-1)/5}$. The argument applies to every maximal acyclic
matching, and therefore to one of minimum cardinality.
\end{proof}

Because every cubic graph has even order, the two parity cases are numerically
equivalent to the unified inequality $n\leq 5k+1$: when $k$ is even, the
right-hand side is odd and integrality sharpens it to $n\leq 5k$. The split
remains structurally informative, since the extra vertex is possible only for
odd $k$, as the example $n=16,k=3$ demonstrates.

\section{Equality constructions for every even order}
\label{sec:constructions}

Throughout this section, $P_r$ and $C_r$ denote the path and cycle on $r$
vertices. For $P_{2s}=v_1\ldots v_{2s}$, its alternating matching is
\[
  \{v_1v_2,v_3v_4,\ldots,v_{2s-1}v_{2s}\}.
\]
For each even $n\geq 4$, we construct a connected simple cubic graph of order
$n$ together with a maximal acyclic matching of size
$\ceil{(n-1)/5}$. By \Cref{thm:sharp-lower}, every displayed matching is
automatically minimum.

\subsection{Base construction for odd matching size}

Let $k_0=2t+1$ be odd with $t\geq 1$. Begin with $k_0$ vertex-disjoint copies
of $K_2$ in $S$, and include all their edges in $M$. Designate one copy as
the root component $R$. A remaining degree slot of a vertex means an
unassigned incident edge needed to bring its degree to three.

All attachments below follow one canonical pairing, so the construction is a
single well-defined graph rather than a family of bijections. For a matched
path $P_{2s}=v_1\ldots v_{2s}$, list remaining slots in vertex order along the
path, repeating each endpoint once for every unused incidence: $v_1$, $v_1$,
then each internal vertex once, then $v_{2s}$, $v_{2s}$. In particular a
copy of $K_2$ contributes the four-slot list $v_1,v_1,v_2,v_2$. For a
$U$-path $x_1\ldots x_r$ or $U$-cycle $z_1\ldots z_r$, list vertices in that
path or cyclic order from the designated start. Pair the two lists in this
order. New $K_2$ components are introduced in the order they are first used.

\begin{enumerate}[leftmargin=2em]
  \item Place a path $P_4=x_1x_2x_3x_4$ in $U$ and attach $x_1,x_2,x_3,x_4$
  to the four canonical slots of $R$ in that order. Join $x_1$ and $x_4$ to
  two distinct unused $K_2$ components $R_1$ and $R_4$, using one slot of
  each.

  \item If $t=1$, declare both $R_1$ and $R_4$ terminal. If $t>1$, declare
  $R_1$ terminal and $R_4$ active. At each subsequent step, take the
  canonical slot list of the active $K_2$, use its first slot for the already
  present upward edge to the previous $U$-path, and pair the remaining three
  slots with a new path $P_3=y_1y_2y_3$. Join $y_1$ to a new unused $K_2$
  (declared terminal) and $y_3$ to a new unused $K_2$ (declared the next
  active component). Repeat $t-1$ times, then declare the final active
  component terminal.

  \item Process terminal $K_2$ components in the order they were designated.
  Each has three remaining canonical slots; attach a triangle
  $z_1z_2z_3$ by pairing those slots with $z_1,z_2,z_3$ in order.
\end{enumerate}

The procedure uses $1+2+2(t-1)=2t+1=k_0$ copies of $K_2$ and leaves
$t+1$ of them terminal. Every degree slot is used once, no loop or parallel
edge is created, every vertex has degree three, and the attachment sequence
keeps the graph connected.

The graph $G[S]$ is the disjoint union of $k_0$ copies of $K_2$, so $M$ is
acyclic. The endpoints of every edge of the root $P_4$ meet $R$; all vertices
of each $P_3$ meet its active $K_2$; and all vertices of each triangle meet
one terminal $K_2$. \Cref{lem:blocking} therefore blocks every edge of
$G[U]$, so $M$ is maximal. Moreover,
\[
  |U|=4+3(t-1)+3(t+1)=6t+4=3k_0+1,
  \qquad
  n=2k_0+|U|=5k_0+1.
\]

\subsection{Residue-class adjustment}

The replaced leaf is the first designated terminal, namely the $K_2$ joined
to $x_1$. In the base construction that $K_2$ is joined to the rest of the
graph by a single edge, and the other end of that edge is a vertex $w\in U$.

For an integer $q\geq 0$, replace this terminal $K_2$ by the alternating path
$P_{2(q+1)}=v_1\ldots v_{2q+2}$ in $S$ and replace the triangle by $C_{2q+3}$
in $U$. Include the $q+1$ alternating path edges in $M$. Form the canonical
slot list of the path, which has $2q+4$ entries. Reserve its first entry
for the reattachment edge $v_1w$, and pair the remaining $2q+3$ slots with
the cycle vertices $z_1,\ldots,z_{2q+3}$ in cyclic order.

Because $w\in U$, the reattachment edge is a cut edge between $S$ and $U$ and
not an edge inside $S$; this is what keeps $G[S]$ a forest. Indeed the
replacement piece contributes the single new $F$-component $P_{2(q+1)}$, which
is a tree, and it is joined to the rest of $S$ by no edge at all. The degree
of $w$ is unchanged, since the edge it previously received from the terminal
$K_2$ is replaced by $v_1w$. The operation therefore preserves simplicity,
connectivity, cubicity, and the fact that $G[S]$ is a forest. Both endpoints
of each new cycle edge meet the same $S$-path component, so \Cref{lem:blocking}
blocks every new $H$-edge and maximality is preserved; the blocking of every
old $H$-edge is untouched, because each such edge is blocked through the root
$K_2$ or an active $K_2$, none of which is modified. Taking $q=0$ recovers the
original leaf piece, so the operation is consistent with the base
construction. \Cref{fig:construction} shows the resulting
component-incidence skeleton for $t=2$ and $q=1$. Relative to $K_2+C_3$, the
replacement adds $q$ matching edges and $4q$ vertices. If $k=k_0+q$, then
\begin{equation}
  n=(5k_0+1)+4q=5k+1-q.
  \label{eq:residue}
\end{equation}

\begin{figure}[t]
\centering
\begin{tikzpicture}[x=1.8cm,y=1.4cm]
  \node[htree]  (X0) at (0,2.4) {$X_0$\\$P_4$};
  \node[htree]  (X1) at (3.7,2.4) {$X_1$\\$P_3$};
  \node[fcomp]  (R)  at (-1,0.8) {$R$\\root $K_2$};
  \node[fcomp]  (L)  at (0,0.8) {$L=P_4$\\replacement};
  \node[fcomp]  (A1) at (1,0.8) {$A_1$\\active $K_2$};
  \node[fcomp]  (A2) at (3.3,0.8) {$A_2$\\terminal $K_2$};
  \node[fcomp]  (B2) at (4.7,0.8) {$B_2$\\terminal $K_2$};
  \node[hcycle] (Z)  at (0,-0.8) {$Z$\\$C_5$};
  \node[hcycle] (ZA) at (3.3,-0.8) {$Z_A$\\$C_3$};
  \node[hcycle] (ZB) at (4.7,-0.8) {$Z_B$\\$C_3$};
  \draw (X0)--(R) (X0)--(L) (X0)--(A1);
  \draw (X1)--(A1) (X1)--(A2) (X1)--(B2);
  \draw (L)--(Z) (A2)--(ZA) (B2)--(ZB);
  \node[font=\scriptsize] at (4.65,2.9) {$t=2,\ q=1$};
\end{tikzpicture}
\caption{Component-incidence skeleton for $t=2,q=1$. Circles are tree
components of $H=G[U]$, rectangles are components of $F=G[S]$, and triangles
are cycle components of $H$.}
\label{fig:construction}
\end{figure}

\subsection{Covering every even order}

Fix an even integer $n\geq 4$, set
\[
  k=\ceil{\frac{n-1}{5}},
  \qquad
  q=5k+1-n.
\]
The definition of the ceiling gives $k-1<(n-1)/5\leq k$, hence
$5k-4<n\leq 5k+1$ and $q\in\{0,1,2,3,4\}$. Put $k_0=k-q$. Since
$k_0=n-4k-1$ and both $n$ and $4k$ are even, $k_0$ is odd.

Only three cases occur. As $q\leq 4$ we have $k_0\geq k-4$, and $k\geq 1$ for
every even $n\geq 4$. If $k_0\leq -5$ then $q=k-k_0\geq 6$, which is
impossible; if $k_0=-3$ then $q=k+3\leq 4$ forces $k=1$ and hence
$n=5k+1-q=2$, excluded by $n\geq 4$. Therefore $k_0\in\{-1,1\}$ or
$k_0\geq 3$, and these three cases are treated in turn.

If $k_0\geq 3$, use the base construction and the size-$q$ leaf replacement.
Equation \eqref{eq:residue} gives exactly $n$ vertices and $k$ matching edges.

If $k_0=1$, then $q=k-1\leq 4$, so $k\leq 5$ and
$n=5k+1-(k-1)=4k+2$, that is, $n\in\{6,10,14,18,22\}$. Take
$F=P_{2k}$ with its alternating matching and $H=C_{2k+2}$, and pair the
$2k+2$ canonical slots of $F$ with the cycle vertices in cyclic order. This graph
has $4k+2=n$ vertices, and every $H$-edge is jointly blocked by the unique
$F$-component. This family is not the $k_0=1$ instance of
\Cref{eq:residue}: with no skeleton to reattach, the path keeps all
$2k+2$ of its slots instead of reserving one, so the attached cycle is
$C_{2k+2}$ rather than $C_{2k+1}$. The distinction is forced, since a piece
$P_{2k}+C_{2k+1}$ would have the odd order $4k+1$ and no cubic graph has odd
order.

It remains to treat $k_0=-1$, where $q=k+1\leq 4$ gives $k\leq 3$ and
$n=5k+1-(k+1)=4k$, that is, $n\in\{4,8,12\}$. For $K_4$, any edge is a maximal
acyclic matching. For the cube $Q_3$, label vertices by three-bit strings and
take
\[
  M=\{000\text{--}001,\ 110\text{--}111\}.
\]
For the hexagonal prism, let the two 6-cycles be
$a_0\ldots a_5$ and $b_0\ldots b_5$ with vertical edges $a_ib_i$, and take
\[
  M=\{a_0b_0,a_2b_2,a_4b_4\}.
\]
In each case, direct application of \Cref{lem:blocking} proves maximality.

\begin{theorem}
\label{thm:extremal}
For every even $n\geq 4$, there exists a connected simple cubic graph $G$ of
order $n$ such that
\[
  \acnum(G)=\ceil{\frac{n-1}{5}}.
\]
Consequently,
\[
  f_3(n)=\ceil{\frac{n-1}{5}}
  \qquad(n\geq 4\text{ even}).
\]
\end{theorem}

\begin{proof}
The constructions above provide a maximal acyclic matching of size
$k=\ceil{(n-1)/5}$, so $\acnum(G)\leq k$.
\Cref{thm:sharp-lower} gives $\acnum(G)\geq k$ for every connected cubic graph.
\end{proof}

\begin{remark}
\label{rem:non-uniqueness}
The constructions establish sharpness but do not classify the extremal graphs.
At order 16 the census contains four graphs with $\acnum(G)=3$, and only three
of them arise from the $t=1$ base construction under some assignment of degree
slots. The fourth replaces the root $P_4$ by a single unmatched vertex joined
to three leaf pieces $K_2+C_3$, one slot per piece. Both configurations
realise $c=3$ and $t=1$ and hence the same value of \eqref{eq:master-count};
what differs is which tree component of $H$ carries the incidences. Any
classification of the extremal graphs must therefore account for at least
these two mechanisms.
\end{remark}

\section{Algorithmic consequences: a theory-guided exact solver}
\label{sec:solverc}

This section translates the proved structure into Solver C, an exact algorithm
specialized to nonempty connected simple cubic graphs. Solver C intentionally
uses \Cref{thm:sharp-lower} and its counting consequences. It is therefore a
\emph{post-theorem} algorithm and must not replace Solvers A and B or the
endpoint-set audit in the non-circular verification chain.

\subsection{The edge-interaction forest}

Let $e=ab$ and $f=cd$ be two candidate graph edges.
\begin{definition}
\label{def:interactions}
The pair $(e,f)$ has:
\begin{enumerate}[label=(\roman*),leftmargin=2em]
  \item an \emph{endpoint conflict} if
  $\{a,b\}\cap\{c,d\}\neq\varnothing$;
  \item a \emph{cycle conflict} if the endpoint sets are disjoint and
  \[
    \lambda(e,f)
    :=|E_G(\{a,b\},\{c,d\})|\geq 2;
  \]
  \item a \emph{link} if the endpoint sets are disjoint and
  $\lambda(e,f)=1$.
\end{enumerate}
For a set $M$ of pairwise endpoint-disjoint edges, the \emph{link graph}
$L(M)$ has vertex set $M$ and an edge $ef$ for each linked pair.
\end{definition}

\begin{proposition}[interaction-forest characterization]
\label{prop:interaction-forest}
A set $M\subseteq E(G)$ is an acyclic matching if and only if:
\begin{enumerate}[label=(\roman*),leftmargin=2em]
  \item no selected pair has an endpoint conflict;
  \item no selected pair has a cycle conflict; and
  \item $L(M)$ is a forest.
\end{enumerate}
\end{proposition}

\begin{proof}
The endpoint-conflict condition is precisely the requirement that $M$ be a
matching, so assume $M$ is a matching and contract every edge of $M$ in the
induced graph $G[V(M)]$. Every vertex of $G[V(M)]$ is covered by $M$, so the
contraction sends each pair of endpoints to a single vertex, and every
remaining edge of $G[V(M)]$ joins two distinct such vertices.

Suppose first that $M$ has a cycle conflict, say $\lambda(e,f)\geq 2$ for
$e=ab$ and $f=cd$. Two distinct edges of $E_G(\{a,b\},\{c,d\})$ together with
$e$ and $f$ contain a cycle of $G[V(M)]$ of length three or four, so
$G[V(M)]$ is not a forest. Suppose next that $M$ has no cycle conflict but
$L(M)$ contains a cycle $f_1f_2\ldots f_sf_1$. Choosing the unique cross edge
realising each link and traversing the matching edges $f_i$ in between yields
a cycle of $G[V(M)]$, so again $G[V(M)]$ is not a forest.

Conversely, suppose $G[V(M)]$ contains a cycle $Q$. Contracting maps $Q$ to a
closed walk $W$ in $L(M)$ of positive length, since $Q$ is not contained in a
single matched pair. If $W$ immediately backtracked, $Q$ would use two
distinct edges of $E_G(\{a,b\},\{c,d\})$ for one pair, that is, a cycle
conflict. In the absence of a cycle conflict, $W$ is a closed walk without
immediate backtracking and therefore contains a cycle of $L(M)$. Hence
$G[V(M)]$ is a forest exactly when there is no cycle conflict and $L(M)$ is a
forest.
\end{proof}

The characterization also gives an exact maximality test. Suppose $M$ is
acyclic and $e\notin M$ has no endpoint or cycle conflict with a selected edge.
Adding $e$ inserts one new vertex into $L(M)$, adjacent to all selected edges
linked to $e$. This preserves the forest exactly when those neighbors lie in
pairwise distinct components of $L(M)$. Therefore $M$ is maximal if and only
if every unselected edge either has a hard conflict (endpoint or cycle), or
has two link-neighbors already lying in the same link-forest component.

\subsection{Rollback-DSU search}

Solver C precomputes all endpoint conflicts, cycle conflicts, and links in
$\bigO(m^2)$ time. During a branch, each selected graph edge is an active
vertex of the link forest. A rollback disjoint-set union structure maintains
its components. To add a candidate $e$, the algorithm checks whether two
current link-neighbors of $e$ already have the same root. If so, adding $e$
closes a cycle; otherwise $e$ is activated and joined to those distinct roots.
A snapshot permits constant-stack rollback after the recursive call.

Search is performed in increasing target cardinality. Capacity bounds use the
number of available suffix edges and their endpoint union. Near a leaf of the
search tree, Solver C also applies a rigid-maximality rule. If an unselected
edge can have at most one link-neighbor in every completion of the current
branch, it can never be cycle-blocked. It must therefore be selected itself or
hard-conflicted by a future selected edge. These necessary alternatives form
small hitting-set constraints; the implementation invokes the rule only when
at most two matching slots remain and solves the resulting hitting-set check
exactly.

\begin{algorithm}[t]
\caption{Solver C: theorem-guided interaction-forest search}
\label{alg:solver-c}
\begin{algorithmic}[1]
\Require A nonempty connected simple cubic graph $G$
\State Precompute endpoint conflicts, cycle conflicts, and links
\State $\ell\gets\ceil{(|V(G)|-1)/5}$
\For{$k=\ell,\ell+1,\ldots,\floor{|V(G)|/2}$}
  \State $D\gets\min\{k-1,\,5k+1-|V(G)|\}$
  \If{$D<0$}
    \State \textbf{continue}
  \ElsIf{$D=0$}
    \State $M\gets\Call{ZeroDefectDomination}{G,k}$
  \Else
    \State $M\gets\Call{InteractionForestSearch}{G,k,D}$
  \EndIf
  \If{$M\neq\text{failure}$}
    \State Verify $M$ with the definition-level matching/forest/maximality routine
    \State \Return $(k,M)$
  \EndIf
\EndFor
\end{algorithmic}
\end{algorithm}

\subsection{The theorem-derived defect budget}

For an acyclic $k$-edge matching, let $\delta=|E(L(M))|$. Since $L(M)$ is a
forest with $k$ vertices, its component count is
\[
  c=k-\delta.
\]
The variables $c$ and $t$ are the same component counts used in
\eqref{eq:master-count}. From
\[
  n=4k+2(c-t)
  \quad\text{and}\quad
  c\leq 2t+1,
\]
we obtain $n\leq 4k+c+1$, hence $c\geq n-4k-1$. Therefore
\[
  \delta=k-c\leq 5k+1-n.
\]
Also $c\geq 1$, so $\delta\leq k-1$. Every solution at target $k$ must satisfy
\begin{equation}
  \boxed{\delta\leq\min\{k-1,\,5k+1-n\}.}
  \label{eq:defect-budget}
\end{equation}
At the first theoretically possible target
$k_0=\ceil{(n-1)/5}$, the quantity $5k_0+1-n$ belongs to
$\{0,1,2,3,4\}$. Thus the first target layer permits at most four link edges.
This bound is safe but theorem-dependent; it is the principal reason Solver C
cannot independently verify \Cref{thm:sharp-lower}.

\subsection{The zero-defect residue class}

When $n=5k+1$, \eqref{eq:defect-budget} forces $\delta=0$. The selected edges
then form an independent set in the union of the endpoint-conflict,
cycle-conflict, and link relations: no two of them are related at all. Every
component of $L(M)$ is a single vertex, so an unselected edge can never have
two link-neighbors in one component, and a single link to a selected edge does
not create a cycle. Maximality is therefore achieved only by selecting an edge
itself or by selecting one of its endpoint or cycle hard conflicts.

Consequently the zero-defect target is an independent domination problem with
one relation for incompatibility and a smaller relation for domination.
Solver C branches on an undominated edge $e$ and selects one member of
\[
  \{e\}\cup N_{\mathrm{hard}}(e).
\]
Every feasible solution must hit this set. Fail-first selection, a safe
coverage lower bound, and memoization of failed
$(\textit{available},\textit{dominated},\textit{slots})$ states complete the
specialized search. No polynomial-time claim is made: the general worst-case
search remains exponential.

\begin{proposition}
\label{prop:solver-c-correct}
Solver C returns $\acnum(G)$ and a corresponding maximal acyclic matching for
every nonempty connected simple cubic graph $G$.
\end{proposition}

\begin{proof}
\Cref{prop:interaction-forest} shows that the precomputed relations and the
rollback structure recognize acyclic matchings exactly. Capacity pruning is
necessary for completing the requested number of disjoint edges. The defect
pruning follows from \eqref{eq:defect-budget}. Every rigid-maximality
constraint is necessary because an edge with at most one possible
link-neighbor cannot become cycle-blocked; the exact small hitting-set test
therefore removes no completion. In the zero-defect case, branching on
$\{e\}\cup N_{\mathrm{hard}}(e)$ is complete for every undominated edge.
The coverage bound is a standard lower bound on the remaining number of
dominating choices, and memoization only reuses failed exact states.

At a complete candidate, the interaction criterion tests maximality exactly,
and the implementation rechecks the returned witness with the independent
definition-level routine. The search starts at the proved lower bound and
increases $k$ only after exhausting the previous target. Hence the first
returned witness is optimal.
\end{proof}

\subsection{A linear vertex kernel}

Consider the parameterized decision problem on connected simple cubic graphs:
given $(G,K)$, decide whether $\acnum(G)\leq K$.

\begin{corollary}
\label{cor:kernel}
\LowAcy{} on connected simple cubic graphs is fixed-parameter tractable with
respect to the solution-size parameter $K$.
\end{corollary}

\begin{proof}
If $\acnum(G)\leq K$, then \Cref{thm:sharp-lower} implies $n\leq 5K+1$.
Thus any instance with $n>5K+1$ is a no-instance and may be replaced by a
fixed no-instance, for example $(K_4,0)$. Every remaining instance already
has at most $5K+1$ vertices. This is an identity kernel on yes-instances: the
bound itself discards oversized inputs, and no further reduction rule is
applied. A cubic graph on at most $5K+1$ vertices has at most
$3(5K+1)/2$ edges, so enumerating edge subsets of size at most $K$ and
checking the definition takes $2^{\bigO(K)}\operatorname{poly}(|G|)$ time.
\end{proof}

The same bound is what Solver~C uses as a search floor. It is recorded here
as an FPT consequence of \Cref{thm:sharp-lower}, not as an independent
kernelization algorithm.

\subsection{Census performance at a glance}

Solver~C is a post-theorem search procedure, not independent evidence for
\Cref{thm:sharp-lower}. On the complete census of $4681$ non-isomorphic
connected simple cubic graphs of orders $4,6,\ldots,16$ it agrees with the
theorem-independent bitset solver on every optimum. \Cref{tab:solverc-bench}
records whole-layer median times: at order $16$ the median drops from
$94.18$ to $32.46$ seconds and the number of complete candidates from
$4{,}832{,}310$ to $40{,}158$. The timing protocol, per-graph ratios,
equality-family scaling, and ablations are given in \Cref{app:benchmark}.

\begin{table}[t]
\centering
\caption{Paired full-census comparison. Times are medians of whole-layer
repetitions. Speedup is $T_B/T_C$, and reduction compares complete candidates.}
\label{tab:solverc-bench}
\scriptsize
\resizebox{\textwidth}{!}{%
\begin{tabular}{rrrrrrrrr}
\toprule
$n$ & Graphs & B time (s) & C time (s) & Speedup & B candidates & C candidates & Reduction & Disc.\\
\midrule
4  & 1     & 0.0008 & 0.0010 & $0.88\times$ & 2         & 1      & 50.00\% & 0\\
6  & 2     & 0.0053 & 0.0040 & $1.33\times$ & 5         & 2      & 60.00\% & 0\\
8  & 5     & 0.0127 & 0.0126 & $1.01\times$ & 70        & 5      & 92.86\% & 0\\
10 & 19    & 0.0839 & 0.0896 & $0.94\times$ & 989       & 358    & 63.80\% & 0\\
12 & 85    & 0.4467 & 0.3775 & $1.18\times$ & 11,151    & 179    & 98.39\% & 0\\
14 & 509   & 5.9727 & 4.8901 & $1.22\times$ & 224,027   & 43,691 & 80.50\% & 0\\
16 & 4,060 & 94.1762& 32.4600& $2.90\times$ & 4,832,310 & 40,158 & 99.17\% & 0\\
\midrule
All & 4,681 & 100.6984 & 37.8348 & $2.66\times$ & 5,068,554 & 84,394 & 98.34\% & 0\\
\bottomrule
\end{tabular}%
}
\end{table}

\section{Independent computational verification and reproducibility}
\label{sec:verification}

\subsection{Independent implementations and structural stress tests}

Solvers A and B agree graph by graph on every cubic graph of order at most 12.
The test suite also covers $K_4$, $K_{3,3}$, paths, cycles, non-cubic graphs,
small graph-atlas instances, and a dedicated trap instance that requires
rebuilding the full vertex-induced subgraph after a matching edge is added.

The order-14 and order-16 audit uses a different search space. For each
even-cardinality set $S\subseteq V(G)$, considered in increasing $|S|$, it
checks by disjoint-set union whether $G[S]$ is a forest, recursively searches
for a perfect matching of $G[S]$, and tests every edge $uv$ outside $S$ by
rebuilding $G[S\cup\{u,v\}]$. A $k$-edge acyclic matching has an endpoint set
of size $2k$ satisfying the first two tests; conversely, a perfect matching of
such a set is acyclic. The last test is equivalent to maximality by
\Cref{lem:one-edge}. The first accepted endpoint set therefore certifies the
optimum, not merely a witness value.

\begin{table}[t]
\centering
\caption{Summary of independent implementations and machine checks.}
\label{tab:verification}
\begin{tabularx}{\textwidth}{>{\raggedright\arraybackslash}p{0.33\textwidth}X}
\toprule
Verification item & Scale or result\\
\midrule
Independent graph generators &
GENREG and nauty 2.9.3 \texttt{geng}\\
Canonical graph-set comparison &
$1,2,5,19,85,509,4060$ classes at orders $4$ through $16$; identical
multisets at every order\\
Solver A versus Solver B &
Graph-by-graph agreement through $n=12$\\
Independent optimum audit, $n=14$ &
All 509 graphs: 418 have value 3 and 91 have value 4\\
Independent optimum audit, $n=16$ &
All 4060 graphs: 4 have value 3 and 4056 have value 4\\
Lower-value exclusion &
Every smaller endpoint-set cardinality was exhausted\\
Solver C interaction audit &
Definition-level interaction characterization checked on
$K_4,K_{3,3}$, the cube, Petersen graph, and a 10-vertex prism\\
Solver B versus Solver C &
All 4681 census graphs plus 64 gate instances; zero discrepancies\\
Current package test suite &
12 tests passed\\
\bottomrule
\end{tabularx}
\end{table}

The endpoint-set audit has its own \texttt{graph6} decoder and does not call the
matching enumeration, acyclicity, maximality, or cache routines of Solvers A
or B. Its vertex-set formulation and recursive perfect-matching test avoid
the principal control flow shared by the two edge-based solvers.

Solver C received separate definition-level interaction audits. Across the
five standard cubic graphs in \Cref{tab:verification}, all tested edge subsets
agree on matching feasibility, acyclicity, and maximality. The publication
benchmark adds 64 gate instances and graph-by-graph agreement on the complete
4681-graph census. These checks validate the implementation but, because
Solver C uses the theorem-derived defect budget, they are not independent
evidence for \Cref{thm:sharp-lower}.

\subsection{Environment and reproducibility materials}

The publication performance run used the environment in
\Cref{tab:benchmark-environment}. Both solvers were single-threaded, no other
benchmark jobs ran concurrently, and the same interpreter, graph objects,
timer, and input order were used. The nauty generator was compiled with GCC
under Ubuntu/WSL, while the timed solver calls ran natively on Windows. The
separate endpoint-set proof audit used its previously frozen environment and
intentionally did not import NetworkX.

\begin{table}[t]
\centering
\caption{Environment of the publication Solver B--C benchmark.}
\label{tab:benchmark-environment}
\small
\begin{tabularx}{\textwidth}{>{\raggedright\arraybackslash}p{0.31\textwidth}X}
\toprule
Item & Recorded value\\
\midrule
Operating system & Microsoft Windows 11, build 26200, 64 bit\\
Processor & AMD Ryzen 7 7435H; 8 physical and 16 logical cores\\
Memory & 15.74 GiB RAM\\
Python & CPython 3.11.7, Anaconda 64 bit, MSC v.1916\\
Graph library & NetworkX 3.1 \cite{Hagberg2008}\\
Graph generator & nauty 2.9.3 \texttt{geng}; connected, minimum/maximum degree 3\\
Execution & Single-threaded; \texttt{time.perf\_counter}; no concurrent jobs\\
Repetitions & Five per complete layer for $n\leq12$; three for $n=14,16$\\
Integrity & SHA-256 recorded for Solver B, Solver C, checker, and every census file\\
Benchmark date & 19 August 2026\\
\bottomrule
\end{tabularx}
\end{table}

The public reproducibility repository records the canonical \texttt{geng}
graph collections, order-specific Solver B--C comparison reports, complete
order-14 and order-16 optimum tables, definition-level routines, Solvers A--C,
structural certificates, equality constructors, tests, commands, software
versions, source and census SHA-256 values, preserved benchmark repetition
times, and all 4681 per-graph comparison rows. Artifacts unavailable for
inclusion in the public snapshot are explicitly identified in the repository.
Source code, graph censuses, computational certificates, and per-graph result
tables are publicly available at
\url{https://github.com/Tanghao514/minimum-maximal-acyclic-matchings-cubic}.

The lower bound and equality construction do not depend on the census. A
single explicit graph with a valid matching certificate suffices to refute the
empirical quarter benchmark, while the exact extremal formula follows entirely
from \Cref{sec:blocking,sec:bound,sec:constructions}. Computation discovers
counterexamples, generates structural hypotheses, and supplies
machine-checkable certificates.

\subsection{Threats to validity}

\begin{itemize}[leftmargin=1.5em]
  \item \textbf{Graph-generation independence.}
  GENREG and \texttt{geng} use different generation algorithms and produce
  identical canonical multisets. They share \texttt{labelg} only as a
  post-generation canonicalizer; neither generator uses it to decide which
  graphs to emit.

  \item \textbf{Optimum verification.}
  At orders 14 and 16, the endpoint-set checker exhausts every smaller
  cardinality and uses separate parsing, forest, perfect-matching, and
  extension code. Reproducibility still depends on preserving source,
  commands, inputs, and hashes.

  \item \textbf{Proof dependence of Solver C.}
  Solver C uses \eqref{eq:defect-budget}, which is derived from the main
  theorem. Its agreement with Solver B is an implementation audit and an
  algorithmic consequence, not a non-circular verification of the theorem.

  \item \textbf{Performance scope.}
  The full-census comparison is a paired single-machine benchmark, not a
  portable performance ranking. Only one canonical graph ordering, one Python
  and NetworkX stack, and one hardware platform were tested. Timings below a
  millisecond are particularly sensitive to interpreter and timer noise.

  \item \textbf{Equality-family censoring.}
  Equality orders above 28 use one run, and Solver B timings at orders 46 and
  56 are right-censored at 300 seconds. Their large structured-instance gains
  are reported separately and are not substituted for census averages.

  \item \textbf{Solver scale.}
  The complete census stops at order 16, and no SAT or ILP certificate
  generator has been applied to the full collection at larger orders. The
  proof is independent of that scale limitation.

  \item \textbf{Literature scope.}
  Results for the maximum acyclic matching number do not settle the
  minimum-maximal problem studied here. In particular the bound
  $\maxacnum(G)\geq n/4$ of \cite{FurstRautenbach2018} bounds a parameter that
  dominates $\acnum(G)$, so it neither implies nor contradicts the present
  theorem; the census in \Cref{sec:census} shows that the two parameters do in
  fact separate on cubic graphs.
\end{itemize}

\section{Discussion and open questions}
\label{sec:discussion}

The extremal formula resolves the fixed-order minimum over connected cubic
graphs, and \Cref{cor:kernel} records the resulting fixed-parameter
tractability for the solution-size parameter in this class. It does not settle
the classical per-instance complexity or the behavior under stronger
restrictions.

\begin{itemize}[leftmargin=1.5em]
  \item \textbf{Cubic complexity.}
  What is the classical complexity of deciding $\acnum(G)\leq K$ on connected
  cubic graphs? The known hardness for bipartite graphs of maximum degree six
  and 4-regular graphs does not resolve this class.

  \item \textbf{Exact exponential algorithms.}
  Can the interaction-forest formulation yield a rigorously improved
  worst-case exponential base, beyond the practical pruning used by Solver C?

  \item \textbf{Kernel size.}
  The reduction in \Cref{cor:kernel} is an identity kernel: yes-instances are
  left unchanged. Can a strictly smaller kernel be obtained by genuine
  reduction rules, and are there conditional or unconditional lower bounds on
  the best possible kernel size?

  \item \textbf{Bounded-width algorithms.}
  Maximal acyclic matching feasibility and cardinality minimization are
  expressible in $\mathrm{LinEMSO}_2$, so the meta-theorem of
  \citet{Arnborg1991} gives fixed-parameter tractability for treewidth. Can one
  obtain an explicit single-exponential dynamic program with matching ETH or
  SETH lower bounds?

  \item \textbf{Verifiable computation.}
  Can compact SAT, ILP, or CP-SAT formulations produce independently
  checkable infeasibility certificates and outperform generic edge-subset
  enumeration?

  \item \textbf{Restricted cubic classes.}
  What are the best lower bounds for 3-edge-connected, bipartite, or
  large-girth cubic graphs?

  \item \textbf{Extremal classification.}
  Can all graphs attaining $\ceil{(n-1)/5}$ be characterized, rather than one
  explicit equality family? \Cref{rem:non-uniqueness} shows that at order 16
  the extremal graphs already split into two structural types, distinguished
  by whether the unique tree component of $H$ is the root $P_4$ or an isolated
  vertex. A classification would have to enumerate the admissible
  component-incidence trees realising $c=2t+1$.

  \item \textbf{Higher regular degree.}
  For fixed $r\geq 4$, what are the asymptotic constant and exact remainder
  term of the analogous extremal function over connected $r$-regular graphs?
\end{itemize}

\section{Conclusion}
\label{sec:conclusion}

We determined the exact extremal behavior of minimum maximal acyclic matchings
in connected cubic graphs. Analysis of the unmatched induced subgraph and its
incidence with the forest induced by matching endpoints yields the sharp
parity-sensitive bound, and explicit constructions attain it at every even
order. The complete order-16 census exposes exactly four counterexamples to
the quarter benchmark inherited from the maximum acyclic matching number;
three of them are realised by the equality family and the fourth is not, so
sharpness and classification are genuinely different questions. Independent
graph generation and three theorem-independent exact search routes separate
computational discovery from mathematical proof.

The structural theorem also has direct algorithmic consequences. Contracting
selected matching edges yields an exact edge-interaction forest, the counting
proof supplies a defect budget of at most four at the first feasible target,
and deciding $\acnum(G)\leq K$ on connected cubic graphs is fixed-parameter
tractable because every yes-instance has at most $5K+1$ vertices. These
consequences strengthen the computational-combinatorics aspect of the result
without changing the central conclusion:
\[
  f_3(n)=\ceil{\frac{n-1}{5}}
  \qquad\text{for every even }n\geq 4.
\]

\appendix
\section{Solver C benchmark details}
\label{app:benchmark}

We compared the theorem-independent bitset solver (Solver~B) with Solver~C on
the complete census of $4681$ non-isomorphic connected simple cubic graphs of
orders $4,6,\ldots,16$. The census was generated with nauty 2.9.3 and its
layer sizes were checked as $1,2,5,19,85,509,$ and $4060$. Before timing, the
implementations were required to agree on $64$ gate instances: ten standard
or prism graphs, the four order-$16$ counterexamples, and $50$ fixed-seed
random connected cubic graphs. Every returned witness was also checked by the
definition-level maximal-acyclic-matching routine. The gate and the complete
census had zero discrepancies.

Graphs were decoded before each timed stage. Each solver received the same
in-memory graphs in the same canonical \texttt{graph6} order. After an untimed
warm-up, every complete layer through order $12$ was repeated five times and
the order-$14$ and order-$16$ layers were repeated three times. Solver order
was alternated between repetitions, garbage collection was disabled only
within individual solver calls, and \texttt{time.perf\_counter} measured
solver time. \Cref{tab:solverc-bench} reports medians of whole-layer totals.
The search columns use the number of complete candidates examined by each
implementation; Solver~C recursion nodes are retained separately in the
archived per-graph data and are not substituted for this common leaf-level
measure.

The whole-census aggregate in \Cref{tab:solverc-bench} is a sum of the seven
layer medians. Per-graph ratios answer a different question: across all $4681$
graphs, their median and geometric mean are $4.41\times$ and $3.81\times$.
At order $16$ they are $4.58\times$ and $4.46\times$, with an interquartile
range of $4.00$--$5.11\times$ and a range of $1.08$--$8.22\times$. Solver~C
is faster on every order-$16$ graph. It is slower on $127$ smaller graphs,
all of order at most $14$, where interpreter, preprocessing, and pruning
overhead are large relative to the underlying search.

Grouping within an order is consistent with the intended role of the defect
budget. At order $10$ the geometric-mean speedups for defect budgets one and
at least five are $1.09\times$ and $0.70\times$; at order $14$ the
corresponding figures for budgets two and at least five are $1.44\times$ and
$1.14\times$. Only four order-$16$ census graphs have zero defect, so their
$5.23\times$ geometric-mean speedup is suggestive rather than a basis for a
broad claim.

The complete census is the primary performance experiment. Separately, we
tested the equality construction at orders
$16,20,24,26,28,32,36,46,$ and $56$ with a common $300$-second timeout.
Orders through $28$ used three repetitions and report medians; larger orders
used one run. \Cref{tab:equality-performance} records these structured
scalability results.

\begin{table}[t]
\centering
\caption{Equality-family comparison under a common $300$-second timeout. The
column $D$ is the defect budget $\min\{k-1,5k+1-n\}$ of
\eqref{eq:defect-budget}; rows with $D=0$ use the specialised zero-defect
routine. Ratios above $100$ are reported only as magnitudes. A dash denotes a
speedup censored by a Solver~B timeout.}
\label{tab:equality-performance}
\small
\begin{tabular}{rrrrrrl}
\toprule
$n$ & $\acnum(G)$ & $D$ & B time (s) & C time (s) & Speedup & Status\\
\midrule
16 & 3  & 0 & 0.0143   & 0.00525 & $2.7\times$      & both solved\\
20 & 4  & 1 & 0.0965   & 0.02424 & $4.0\times$      & both solved\\
24 & 5  & 2 & 0.8888   & 0.15152 & $5.9\times$      & both solved\\
26 & 5  & 0 & 1.5255   & 0.01000 & ${>}10^{2}\times$ & both solved\\
28 & 6  & 3 & 8.4682   & 0.96229 & $8.8\times$      & both solved\\
32 & 7  & 4 & 89.2176  & 6.44114 & $13.9\times$     & both solved\\
36 & 7  & 0 & 290.5094 & 0.01720 & ${>}10^{4}\times$ & both solved\\
46 & 9  & 0 & timeout  & 0.02144 & --               & B timeout\\
56 & 11 & 0 & timeout  & 0.02971 & --               & B timeout\\
\bottomrule
\end{tabular}
\end{table}

The five rows with $D=0$ are exactly the orders with $n=5k+1$, namely
$n\equiv 6\pmod{10}$, where the budget of \eqref{eq:defect-budget} collapses
and Solver~C enters the specialised domination routine. On those rows
Solver~C takes between $0.005$ and $0.030$ seconds across the whole range
$16\leq n\leq 56$, whereas Solver~B grows from $0.014$ seconds to a timeout.
The extreme ratios therefore track the blow-up of the unguided baseline rather
than any variation in Solver~C, which is also why the same zero-defect mode
yields $2.7\times$ at order $16$ and more than $10^4\times$ at order $36$.
These are structured members of the equality family, not a random sample, and
the ratios must not be read as average census speedups.

Finally, five Solver~C variants were run three times on a deterministic,
stratified sample of $100$ order-$16$ graphs. On the four zero-defect graphs,
replacing the specialized zero-defect routine by the generic search increased
median search time by approximately $9.7\times$, while disabling the
theorem-derived budget increased it by approximately $15.2\times$. Failed-state
memoization had no measurable benefit on these four found-solution searches.
On the remaining $96$ graphs, removing near-leaf maximality pruning increased
the total recursion-node count from $1853$ to $3268$, although the additional
pruning test did not improve median microbenchmark time. These ablations
identify mechanisms, not portable worst-case guarantees.

\begin{thebibliography}{99}

\bibitem[Arnborg et al.(1991)Arnborg, Lagergren, and Seese]{Arnborg1991}
Arnborg S, Lagergren J, Seese D (1991) Easy problems for tree-decomposable
graphs. Journal of Algorithms 12:308--340.
\url{https://doi.org/10.1016/0196-6774(91)90006-K}

\bibitem[Cameron(1989)]{Cameron1989}
Cameron K (1989) Induced matchings. Discrete Applied Mathematics 24:97--102.
\url{https://doi.org/10.1016/0166-218X(92)90275-F}

\bibitem[Chaudhary and Panda(2021)]{ChaudharyPanda2021}
Chaudhary J, Panda BS (2021) On the complexity of minimum maximal uniquely
restricted matching. Theoretical Computer Science 882:15--28.
\url{https://doi.org/10.1016/j.tcs.2021.05.036}

\bibitem[Chaudhary and Zehavi(2025)]{Zehavi2025}
Chaudhary J, Zehavi M (2025) $p$-matchings parameterized by treewidth.
SIAM Journal on Discrete Mathematics 39:1280--1311.
\url{https://doi.org/10.1137/23M160013X}

\bibitem[Chaudhary et al.(2024)Chaudhary, Mishra, and Panda]{Chaudhary2024}
Chaudhary J, Mishra S, Panda BS (2024) On the complexity of minimum maximal
acyclic matchings. Journal of Combinatorial Optimization 48:10.
\url{https://doi.org/10.1007/s10878-024-01200-3}

\bibitem[Chaudhary et al.(2025)Chaudhary, Mishra, and Panda]{Chaudhary2025}
Chaudhary J, Mishra S, Panda BS (2025) Minimum maximal acyclic matching in
proper interval graphs. Discrete Applied Mathematics 360:414--427.
\url{https://doi.org/10.1016/j.dam.2024.10.012}

\bibitem[Edmonds(1965)]{Edmonds1965}
Edmonds J (1965) Paths, trees, and flowers. Canadian Journal of Mathematics
17:449--467. \url{https://doi.org/10.4153/CJM-1965-045-4}

\bibitem[F{\"u}rst and Rautenbach(2018)]{FurstRautenbach2018}
F{\"u}rst M, Rautenbach D (2018) A lower bound on the acyclic matching number
of subcubic graphs. Discrete Mathematics 341(8):2353--2358.
\url{https://doi.org/10.1016/j.disc.2018.05.010}

\bibitem[F{\"u}rst and Rautenbach(2019)]{Furst2019}
F{\"u}rst M, Rautenbach D (2019) On some hard and some tractable cases of the
maximum acyclic matching problem. Annals of Operations Research 279:291--300.
\url{https://doi.org/10.1007/s10479-019-03311-1}

\bibitem[Goddard et al.(2005)Goddard, Hedetniemi, Hedetniemi, and Laskar]{Goddard2005}
Goddard W, Hedetniemi SM, Hedetniemi ST, Laskar R (2005) Generalized
subgraph-restricted matchings in graphs. Discrete Mathematics 293:129--138.
\url{https://doi.org/10.1016/j.disc.2004.08.027}

\bibitem[Golumbic et al.(2001)Golumbic, Hirst, and Lewenstein]{Golumbic2001}
Golumbic MC, Hirst T, Lewenstein M (2001) Uniquely restricted matchings.
Algorithmica 31:139--154. \url{https://doi.org/10.1007/s00453-001-0004-z}

\bibitem[Hagberg et al.(2008)Hagberg, Schult, and Swart]{Hagberg2008}
Hagberg AA, Schult DA, Swart PJ (2008) Exploring network structure, dynamics,
and function using NetworkX. In: Proceedings of the 7th Python in Science
Conference, pp 11--15

\bibitem[Lov{\'a}sz and Plummer(2009)]{LovaszPlummer2009}
Lov{\'a}sz L, Plummer MD (2009) Matching Theory. American Mathematical
Society, Providence

\bibitem[McKay and Piperno(2014)]{McKayPiperno2014}
McKay BD, Piperno A (2014) Practical graph isomorphism, II. Journal of
Symbolic Computation 60:94--112. \url{https://doi.org/10.1016/j.jsc.2013.09.003}

\bibitem[Meringer(1999)]{Meringer1999}
Meringer M (1999) Fast generation of regular graphs and construction of cages.
Journal of Graph Theory 30:137--146.
\url{https://doi.org/10.1002/(SICI)1097-0118(199902)30:2<137::AID-JGT7>3.0.CO;2-G}

\bibitem[Panda and Chaudhary(2023)]{Panda2023}
Panda BS, Chaudhary J (2023) Acyclic matching in some subclasses of graphs.
Theoretical Computer Science 943:36--49.
\url{https://doi.org/10.1016/j.tcs.2022.12.008}

\bibitem[Yannakakis and Gavril(1980)]{Yannakakis1980}
Yannakakis M, Gavril F (1980) Edge dominating sets in graphs. SIAM Journal on
Applied Mathematics 38:364--372. \url{https://doi.org/10.1137/0138030}

\end{thebibliography}

\section*{Statements and Declarations}

\paragraph{Funding.}
The author declares that no funds, grants, or other support were received
during the preparation of this manuscript.

\paragraph{Competing interests.}
The author has no relevant financial or non-financial interests to disclose.

\paragraph{Author contributions.}
Hao Tang conceived the study, implemented the algorithms, designed and
conducted the experiments, performed the structural analysis, developed the
proofs, and wrote the manuscript.

\paragraph{Data and code availability.}
Source code, graph censuses, computational certificates, and per-graph result
tables are publicly available at
\url{https://github.com/Tanghao514/minimum-maximal-acyclic-matchings-cubic}.

\paragraph{Use of AI-assisted tools.}
AI-assisted tools were used for exploratory prompts, code review, and language
polishing. They were not listed as authors and did not assume responsibility
for mathematical or computational claims. The author takes full
responsibility for the final manuscript.

\end{document}
